\documentclass[11pt]{article}

\usepackage[a4paper,margin=1in]{geometry}
\usepackage[T1]{fontenc}
\usepackage[utf8]{inputenc}
\usepackage{lmodern}
\usepackage{graphicx}
\usepackage{hyperref}

\hypersetup{
  colorlinks=true,
  linkcolor=black,
  urlcolor=blue,
  pdftitle={Oracle Engine v2 and Deterministic Routing},
  pdfauthor={OpenAI Codex},
  pdfsubject={Stability Oracle architectural update}
}
\setcounter{secnumdepth}{3}
\setcounter{tocdepth}{2}

\title{Oracle Engine v2 And Deterministic Routing\\Project Paper}
\author{Stability Oracle Architectural Update}
\date{March 24, 2026}

\begin{document}
\maketitle
\tableofcontents
\newpage

\section{Purpose}

This paper documents the next additive architectural step in the Stability Oracle project:

\begin{itemize}
\item Oracle Engine v2
\item deterministic perturbation evaluation
\item deterministic trace payloads
\item explicit route selection for future oracle domains
\end{itemize}

This is an engineering paper, not a theory paper.
It does not expand metric semantics, ontology, or physical claims.
Its purpose is to record how the repository now turns the existing metric and policy layers into a reusable decision surface and how it safely routes future oracle domains without pretending they are implemented.

\section{Background}

The repository already exposed a bounded structural survivability surface:

\begin{itemize}
\item invariant state contract
\item perturbation contract
\item pluggable metric registry
\item deterministic classification policy
\item CLI and agent-facing skill interface
\end{itemize}

That system answered one question:

\begin{quote}
Given a state and a perturbation, does meaningful structure remain recoverably coherent?
\end{quote}

The external conceptual reference in \texttt{Images/The\_Oracle\_Paradigm.pdf} suggests a broader stack with multiple oracle domains.
This update does not implement that broader stack.
It only introduces the safe controller surface required to route into it later.

\section{Current Layering}

The repository now contains two distinct but compatible evaluation surfaces.

\subsection{Application Adapter Surface}

The original application skill remains intact.
It calls an external HAOS oracle CLI and converts the returned telemetry into the bounded public report format.
This preserves compatibility with the documented CLI and external integration examples.

\subsection{Local Logical Surface}

Oracle Engine v2 adds a fully local logical layer that operates on typed \texttt{State} and \texttt{Perturbation} objects.
It uses the in-repository metric registry and classifier policy directly.

\begin{figure}[h]
  \centering
  \includegraphics[width=\textwidth]{../../Images/Stability Oracle Evaluation Layer Overview.png}
  \caption{Existing Stability Oracle evaluation layer overview.}
\end{figure}

\section{Oracle Engine v2}

\subsection{Goal}

Oracle Engine v2 exists to turn the metric and policy layers into a reusable local decision engine.
It supports three deterministic operations:

\begin{itemize}
\item transition evaluation
\item perturbation evaluation
\item ordered batch scanning
\end{itemize}

\subsection{Engine Interface}

The engine exposes:

\begin{verbatim}
engine.evaluate_transition(before, after)
engine.evaluate(before, perturbation)
engine.scan(before, perturbations)
\end{verbatim}

These operations are side-effect free with respect to their inputs.
They do not mutate state objects, they do not invoke learned models, and they do not depend on external repositories.

\subsection{Perturbation Engine}

The local perturbation engine applies a fixed deterministic transform order:

\begin{enumerate}
\item node drop
\item cluster split
\item edge drop
\item conservative noise passthrough
\end{enumerate}

The noise field is intentionally conservative in this version.
It is accepted by the contract, validated, and preserved as part of the perturbation language, but it does not invent new structure-changing semantics until a deterministic and explicit rule is defined.

\subsection{Trace Payload}

Oracle Engine v2 extends \texttt{OracleResult} with an optional deterministic trace payload.
The trace remains compact and JSON-friendly.
It includes:

\begin{itemize}
\item input node count
\item input edge count
\item output node count
\item output edge count
\item normalized metric vector
\item coherence score
\item spread
\item floor trigger status
\item policy version
\end{itemize}

This makes local execution inspectable without bloating the main result contract.

\section{Deterministic Reporting}

The rule-based reporting layer turns classification results into short bounded explanations.
These explanations are not generated by an LLM.
They are policy-coupled summaries such as:

\begin{itemize}
\item stable because all normalized dimensions stayed above threshold and spread remained bounded
\item marginal because the survivability floor held but coherence was uneven
\item unstable because at least one core dimension fell below the survivability floor
\end{itemize}

This preserves interpretability and reproducibility.

\section{Routing Layer}

\subsection{Motivation}

The project now distinguishes three possible oracle route labels:

\begin{itemize}
\item \texttt{PHYSICAL}
\item \texttt{LOGICAL}
\item \texttt{FOUNDATIONAL}
\end{itemize}

Only \texttt{LOGICAL} is implemented.
The routing layer does not pretend otherwise.

\subsection{Routing Context}

The new routing context is a typed JSON-friendly contract containing bounded fields such as:

\begin{itemize}
\item candidate identifier
\item domain hint
\item uncertainty level
\item explicit boolean triggers for empirical, structural, or foundational evaluation
\item prior route
\item optional JSON notes
\end{itemize}

This context is validated and normalized.
No hidden fields are interpreted.

\subsection{Route Policy}

The route policy is explicit and rule-based.
Its priority order is:

\begin{enumerate}
\item foundational requirement
\item structural stability requirement
\item empirical search requirement
\item known domain hint
\item unknown fallback
\end{enumerate}

Confidence is intentionally simple:

\begin{itemize}
\item \texttt{1.0} for explicit boolean triggers
\item \texttt{0.8} for direct domain hint matches
\item \texttt{0.4} for unknown fallback
\end{itemize}

The router therefore remains interpretable and safe to inspect.

\subsection{Dispatch Semantics}

The router exposes two operations:

\begin{itemize}
\item \texttt{decide(context)}: select route only
\item \texttt{dispatch(context, before, after, perturbation)}: select and execute when safe
\end{itemize}

Dispatch behavior:

\begin{itemize}
\item \texttt{LOGICAL}: execute Oracle Engine v2 when a logical engine is available
\item \texttt{PHYSICAL}: return bounded stub
\item \texttt{FOUNDATIONAL}: return bounded stub
\item \texttt{UNKNOWN}: return bounded stub
\end{itemize}

Unimplemented routes therefore fail closed without pretending to evaluate.

\section{Bounded Stub Design}

The stub response format is intentionally small:

\begin{verbatim}
{
  "status": "stub",
  "route": "PHYSICAL",
  "message": "...",
  "implemented": false
}
\end{verbatim}

This keeps route selection usable today without creating false domain logic.

\section{Compatibility}

\subsection{CLI}

The documented CLI remains compatible.
Legacy payloads still use the external adapter-backed skill surface.
New local-engine payload shapes are detected conservatively:

\begin{itemize}
\item \texttt{before + after}
\item \texttt{before + perturbation}
\item \texttt{before + perturbations}
\end{itemize}

No documented command was removed.

\subsection{Public Dataclasses}

Existing dataclasses remain in place.
The only additive extension is the optional \texttt{trace} field on \texttt{OracleResult}.

\section{Testing}

The repository test suite now covers:

\begin{itemize}
\item perturbation immutability
\item deterministic transition evaluation
\item equivalence between derived and explicit transitions
\item ordered batch scan behavior
\item explicit engine exceptions
\item deterministic explanation output
\item CLI compatibility
\item routing priorities
\item stub behavior for unimplemented routes
\item logical route dispatch
\end{itemize}

At the time of this paper, the standalone suite passes end-to-end.

\section{What This Update Does Not Do}

This update does not implement:

\begin{itemize}
\item physical surrogate search
\item foundational admissibility logic
\item learned routing
\item adaptive policy tuning
\item remote execution
\item ontology expansion
\end{itemize}

The repository now supports those future surfaces only as route labels and bounded placeholders.

\section{Operational Summary}

The project now supports two important patterns:

\begin{verbatim}
from haos_skill.oracle.engine import build_default_engine
engine = build_default_engine()
result = engine.evaluate(before_state, perturbation)
\end{verbatim}

and:

\begin{verbatim}
from haos_skill.routing.router import build_default_router
from haos_skill.routing.context import RoutingContext

router = build_default_router()
decision = router.decide(RoutingContext(requires_structural_stability=True))
\end{verbatim}

This means the repository has moved from a pure metric toolkit to a bounded logical engine with an explicit future-facing route controller.

\section{Conclusion}

The Stability Oracle repository now has:

\begin{itemize}
\item a local deterministic oracle engine
\item deterministic perturbation application
\item bounded traceable results
\item explicit deterministic oracle routing
\item safe stubs for unimplemented domains
\end{itemize}

This is still a small system by design.
That is a strength.
It keeps the current logical surface inspectable while creating a disciplined path toward a larger oracle stack without claiming that larger stack already exists.

\end{document}
